\section*{Chapter \ref{ch:rs:build-a-research-workflow} --- Build a Research Workflow}

\begin{solution}{exo:rs:build-a-research-workflow:1}
The same parameters run through different code give different results; a key without the code's hash would serve the old result after the code changed (the
tear sheet that gained a \omterm{def:rs:performance-measurement:hit}{hit rate} would have been read from the cache without it).
\end{solution}

\begin{solution}{exo:rs:build-a-research-workflow:2}
So that a change upstream that leaves an intermediate output unchanged stops there: the downstream keys depend only on what the inputs contain, not on how they
were produced (early cutoff).
\end{solution}

\begin{solution}{exo:rs:build-a-research-workflow:3}
Snapshot, universe, features, cards and level~1 run, since each reads the snapshot; only the snapshot's output changes, so blend, portfolio and the tear sheet
are read from the cache.
\end{solution}

\begin{solution}{exo:rs:build-a-research-workflow:4}
It enforces keeping track of how each result was produced (1), recording intermediate results (5), noting seeds (6, by flagging unseeded stages in
reproduction), and part of archiving versions (3, the environment record). Avoiding manual steps (2), version control (4), keeping plot data (7), hierarchical
output (8), linking statements to results (9) and public access (10) remain practices, though the manifest makes 9 easy.
\end{solution}

\begin{solution}{exo:rs:build-a-research-workflow:5}
A change to the helper changes the stage's output without changing its key, so the cache serves a stale result. Hash the functions the stage calls (its
module's source, or the transitive closure of the functions it uses), or version helpers explicitly as parameters.
\end{solution}

\begin{solution}{exo:rs:build-a-research-workflow:6}
The environment: library versions (NumPy's summation and BLAS builds change the order of floating-point operations), the number of threads, the platform. The
manifest's environment record says what to compare; bitwise reproduction across environments may require pinned versions and single-threaded reductions (Book~4,
chapter~25).
\end{solution}

\begin{solution}{exo:rs:build-a-research-workflow:7}
Without a seed, \texttt{reproduce} returns the tear sheet alone. A seed derived from the hash of the stage's inputs and parameters is fixed for a given computation,
so it reproduces, and it is chosen by no one, so it cannot be tuned: a result that depends on the seed is still exposed by varying the inputs.
\end{solution}

\begin{solution}{exo:rs:build-a-research-workflow:8}
Version control keeps the code, not the data it read (the vendor revised them), not the environment, not the order in which the cells were run, not the seeds,
nor which version produced the published numbers. A manifest records all of them.
\end{solution}

\begin{solution}{pb:rs:build-a-research-workflow:1}
\begin{enumerate}
\item Snapshot; universe (snapshot); features (snapshot, universe); cards (snapshot, features); blend (features, universe); portfolio (blend, universe); level~1
(portfolio, snapshot); tear sheet (level~1).
\item The key: name, code hash, parameters, seed, the inputs' output hashes. The manifest: for each stage its key, code hash, parameters, seed, dependencies, input
hashes, output hash and whether it ran, and the environment.
\item It reads every stage from the cache: all keys are stored.
\item Momentum's rank IC 0.0076 ($t = 1.39$), the reversal's 0.0379 ($t = 16.05$) over 2\,267 days; the monthly book's Sharpe ratio after costs 0.08 (standard error
0.33, bootstrap interval $-0.67$ to 0.74), turnover 3.5\% a day.
\item Snapshot, universe, features, cards, level~1 re-run; only the snapshot's output changes.
\item All eight re-run, and all eight outputs change.
\item Blend, portfolio, level~1 and tear sheet.
\item The tear sheet only.
\item The same output hash for all eight stages: an empty list of differences.
\item Reproduction flags the tear sheet, and only it.
\item Python and NumPy versions and the platform; when a reproduction fails on another machine, or floating-point results differ across library versions.
\item \textbf{Named result.} A revision outside the universe re-runs five stages and changes one output (the snapshot's); a revision inside it re-runs and changes all
eight; reproducing the manifest from scratch returns identical output hashes for all eight stages.
\item As a trial of its family, with the hash of the manifest's stage keys as its code identity and the snapshot's output hash as its data identity, and its
metrics.
\item Reproduction, \omterm{def:rs:point-in-time-data-and-the-biases:lookahead}{point-in-time data} (partly: the snapshot's identity and revisions), the \omterm{def:rs:the-research-process:log}{trial count} (through registration); not the \omterm{def:rs:vectorised-backtests:fidelity}{fidelity level}, the
costs at size, or whether the text matches the numbers.
\item The book is rebuilt monthly and the reversal decays within a day (chapter~6's horizon mismatch).
\item As a manifest with its cache and environment, registered in the log: one command to re-run it, and a diff to say what the vendor's revision changed.
\item Code hashes that follow called functions, a shared cache, enforced \omterm{def:rs:build-a-research-workflow:repro}{environment locks}, and review sign-off recorded with the manifest.
\item One that its manifest produces again, bit for bit, from the same data, code, parameters, seeds and environment.
\end{enumerate}
\end{solution}

\begin{solution}{iq:rs:build-a-research-workflow:1}
Revised or re-snapshotted data (vendor \omterm{def:rs:point-in-time-data-and-the-biases:vintage}{restatements}, survivorship fixes), changed code or library versions, different parameters, unseeded randomness, a different
universe or calendar, look-ahead fixed since, or a different environment (threads, platform).
\end{solution}

\begin{solution}{iq:rs:build-a-research-workflow:2}
Keep track of how every result was produced; automate every step; version code and data (snapshots); record intermediate results; seed randomness; lock
environments; link every number in a write-up to its source (Sandve and co-authors' ten rules).
\end{solution}

\begin{solution}{iq:rs:build-a-research-workflow:3}
Content-addressed: key each stage by the hash of its code, parameters, seed and its inputs' content hashes; store outputs under their keys; never invalidate
explicitly, since a changed input produces a new key; build keys on output hashes for early cutoff; garbage-collect unreferenced entries.
\end{solution}

\begin{solution}{iq:rs:build-a-research-workflow:4}
Reproduction from a manifest; \omterm{def:rs:point-in-time-data-and-the-biases:lookahead}{point-in-time data} and universe; the \omterm{def:rs:the-research-process:log}{trial count}, deflated Sharpe ratio and PBO; the \omterm{def:rs:vectorised-backtests:fidelity}{fidelity level} of the \omterm{def:rs:vectorised-backtests:backtest}{backtest}; costs and
\omterm{def:rs:capacity-decay-and-crowding:capacity}{capacity} at the intended size; stability over sub-periods; that the text's numbers match the outputs.
\end{solution}

\begin{solution}{iq:rs:build-a-research-workflow:5}
Seed every generator explicitly (never from the clock), use a generator with a documented stream (not a library default that may change), record the seed and the
library version, and make the order of draws independent of threads.
\end{solution}

\begin{solution}{iq:rs:build-a-research-workflow:6}
Find every manifest whose snapshot stage used the vendor's data (the log's data identity); re-run each with the restated snapshot; the diffs say which stages'
outputs changed, and early cutoff says which results are untouched.
\end{solution}
